\section{Benchmark Function Definitions}
\label{app:benchmarks}

\Cref{tab:benchmark-formulas} provides the complete mathematical definitions
of all benchmark functions used in this study.  Each function is a standard
test problem drawn from the optimisation literature.

\begin{table}[H]
\centering
\caption{Mathematical definitions of the six benchmark functions.}
\label{tab:benchmark-formulas}
\renewcommand{\arraystretch}{1.6}
\begin{tabularx}{\textwidth}{@{}lX@{}}
\toprule
\textbf{Function} & \textbf{Definition} \\
\midrule
Sphere
  & $f(\vect{x}) = \displaystyle\sum_{i=1}^{n} x_i^{2}$ \\
Rastrigin
  & $f(\vect{x}) = 10n + \displaystyle\sum_{i=1}^{n}
    \bigl[x_i^{2} - 10\cos(2\pi x_i)\bigr]$ \\
Rosenbrock
  & $f(\vect{x}) = \displaystyle\sum_{i=1}^{n-1}
    \bigl[100(x_{i+1} - x_i^{2})^{2} + (1 - x_i)^{2}\bigr]$ \\
Ackley
  & $f(\vect{x}) = -20\exp\!\Bigl(-0.2\sqrt{\frac{1}{n}
    \sum_{i=1}^{n} x_i^{2}}\Bigr)
    - \exp\!\Bigl(\frac{1}{n}\sum_{i=1}^{n}\cos(2\pi x_i)\Bigr)
    + 20 + e$ \\
Griewank
  & $f(\vect{x}) = \displaystyle\frac{1}{4000}\sum_{i=1}^{n} x_i^{2}
    - \prod_{i=1}^{n}\cos\!\Bigl(\frac{x_i}{\sqrt{i}}\Bigr) + 1$ \\
Schwefel
  & $f(\vect{x}) = 418.9829\,n
    - \displaystyle\sum_{i=1}^{n} x_i\sin\!\bigl(\sqrt{|x_i|}\bigr)$ \\
\bottomrule
\end{tabularx}
\end{table}


\section{Key Source Code Listings}
\label{app:code}

The following listings present the core algorithmic components of the
framework.  The complete source code is available in the accompanying
repository.

\subsection{QPSO Position Update}
\label{app:qpso-update}

The central loop of the QPSO algorithm applies the quantum-inspired
position update to every particle at each iteration.

\begin{minted}{python}
# Core QPSO position update (from qpso.py)
for i, particle in enumerate(self.swarm):
    # Attractor: weighted combination of personal and global best
    phi = r1[i]
    attractor = phi * particle.best_pos + (1 - phi) * self.global_best_pos

    # Quantum displacement
    delta = beta * np.abs(mbest - particle.position) * np.log(1.0 / u[i])

    # Tunnelling: sign determines direction
    particle.position = np.where(
        sign[i] > 0,
        attractor + delta,
        attractor - delta,
    )

    # Enforce search bounds
    particle.position = np.clip(particle.position, self.bounds[0], self.bounds[1])
\end{minted}


\subsection{Weighted Mean Best Computation}
\label{app:wqpso}

The fitness-weighted mean best position biases the swarm reference
point toward higher-quality regions of the search space.

\begin{minted}{python}
# Weighted mean best (WQPSO) — from convergence.py
def compute_weighted_mean_best(
    personal_bests: list[np.ndarray],
    best_scores: list[float],
) -> np.ndarray:
    scores = np.array(best_scores)
    worst = scores.max()
    score_range = worst - scores.min()
    weights = (worst - scores) + score_range * 0.01
    weights /= weights.sum()
    positions = np.array(personal_bests)
    return np.average(positions, axis=0, weights=weights)
\end{minted}


\subsection{Quantum Circuit — Beta-Biased Random Sign}
\label{app:quantum-circuit}

For the \texttt{full} and \texttt{hybrid} modes, the tunnelling sign is
sampled from a single-qubit circuit whose measurement bias is controlled
by the contraction--expansion coefficient~$\beta$.  A qubit initialised in
$|0\rangle$ is passed through a Hadamard gate followed by a parameterised
rotation
\begin{equation}
  R_Y(\theta) = \begin{pmatrix}
    \cos(\theta/2) & -\sin(\theta/2) \\
    \sin(\theta/2) & \phantom{-}\cos(\theta/2)
  \end{pmatrix},
  \label{eq:ry_gate}
\end{equation}
with $\theta = \beta\pi$.  After the combined $H \cdot R_Y(\beta\pi)$
operation, the probability of measuring $|1\rangle$ is
\begin{equation}
  P(|1\rangle) = \sin^2\!\biggl(\frac{\beta\pi}{2}\biggr).
  \label{eq:prob_one}
\end{equation}
When $\beta$ is large (early exploration) $P(|1\rangle) \approx 0.5$, an
approximately unbiased coin flip; as $\beta$ decays toward $0.5$,
$P(|1\rangle) \approx 0.15$, biasing the sign toward $+1$ (outward
displacement) and thereby tightening the effective search radius.  The measured
fraction of $|1\rangle$ outcomes over the shot budget simultaneously provides
the displacement magnitude and, via thresholding at $0.5$, the binary tunnelling
sign.  In \texttt{full} mode, two additional Hadamard-only circuits supply the
cognitive and social weights $r_1$ and $r_2$.  All circuits for a single
iteration are batched into one backend job to minimise communication overhead.

\begin{minted}{python}
# Quantum circuit for beta-biased sign (from quantum_engine.py)
from qiskit import QuantumCircuit

qc = QuantumCircuit(1, 1)
qc.h(0)                     # Hadamard: uniform superposition
qc.ry(beta * np.pi, 0)      # RY rotation: bias toward |1>
qc.measure(0, 0)             # collapse to classical bit

# P(|1>) = sin^2(beta * pi / 2)
# High beta -> P ~ 0.5 (broad exploration)
# Low beta  -> P ~ 0.15 (tight exploitation)
\end{minted}


\subsection{Stagnation Detection and Scatter}
\label{app:stagnation}

When the global best score fails to improve for a configurable number
of consecutive iterations, the worst-performing fraction of the swarm
is re-scattered near the current best position.

\begin{minted}{python}
# Stagnation scatter — from convergence.py
def scatter_worst_particles(particles, global_best_pos, bounds, fraction=0.3):
    ranked = sorted(particles, key=lambda p: p.best_score, reverse=True)
    n_scatter = max(1, int(len(particles) * fraction))
    search_range = bounds[1] - bounds[0]
    spread = search_range * 0.15

    for p in ranked[:n_scatter]:
        noise = np.random.uniform(-spread, spread, size=p.position.shape)
        p.position = np.clip(
            global_best_pos + noise, bounds[0], bounds[1]
        )
        p.best_score = float("inf")
        p.best_pos = p.position.copy()
\end{minted}


\section{Reproduction Guide}
\label{app:reproduction}

To reproduce the experiments reported in this paper, follow these steps:

\begin{enumerate}[leftmargin=*]
    \item \textbf{Clone the repository} and navigate to the project root:
    \begin{minted}{bash}
git clone <repository-url>
cd PY
    \end{minted}

    \item \textbf{Create a virtual environment} and install the package:
    \begin{minted}{bash}
python -m venv .venv
# Linux / macOS:
source .venv/bin/activate
# Windows (PowerShell):
.venv\Scripts\Activate.ps1
# Windows (Git Bash / MSYS2):
source .venv/Scripts/activate

pip install -e ".[dev]"
    \end{minted}

    \item \textbf{Run the quick test} to verify the installation:
    \begin{minted}{bash}
pytest -v
python experiments/run_comparison.py --benchmark sphere --dims 2
    \end{minted}

    \item \textbf{Run the full study} with the publication configuration:
    \begin{minted}{bash}
python experiments/run_full_study.py --config configs/publication.yaml
    \end{minted}
    This produces convergence histories in \texttt{results/raw/},
    figures in \texttt{results/figures/}, and statistical summaries
    in \texttt{results/summary/}.

    \item \textbf{(Optional) Enable quantum modes:}
    \begin{minted}{bash}
pip install -e ".[quantum]"
    \end{minted}
    Then add \texttt{full} or \texttt{hybrid} to the \texttt{qpso\_modes}
    list in your YAML configuration file.
\end{enumerate}
